% Lines 2599-3260 from original attention_transformers_notes_ghost.tex
% Chapter 11: Advanced Architectural Variants

\section{Advanced Architectural Variants}
%==============================================================================

\subsection{Efficient Attention Mechanisms}

The quadratic complexity of attention has motivated numerous efficient variants.

\subsubsection{Linear Attention}

Linear attention approximates softmax attention using kernel methods:

\begin{equation}
\text{LinearAttention}(Q, K, V) = \phi(Q)(\phi(K)^T V)
\end{equation}

where $\phi$ is a feature map. This reduces complexity from $O(n^2d)$ to $O(nd^2)$.

\begin{lstlisting}[language=Python, caption=Linear Attention Implementation]
from __future__ import annotations
from dataclasses import dataclass
from typing import Protocol, runtime_checkable
import torch
import torch.nn as nn
import torch.nn.functional as F
from torch import Tensor

@dataclass(frozen=True, slots=True)
class LinearAttentionConfig:
    d_model: int
    n_heads: int
    feature_map: str = "elu"
    eps: float = 1e-6

@runtime_checkable
class FeatureMap(Protocol):
    def __call__(self, x: Tensor) -> Tensor: ...

class ELUFeatureMap:
    """ELU + 1 feature map for linear attention"""

    def __call__(self, x: Tensor) -> Tensor:
        return F.elu(x) + 1.0

class LinearMultiHeadAttention(nn.Module):
    """Linear complexity multi-head attention using kernel trick"""

    def __init__(self, config: LinearAttentionConfig) -> None:
        super().__init__()
        self.config = config
        self.d_k = config.d_model // config.n_heads

        self.w_q = nn.Linear(config.d_model, config.d_model, bias=False)
        self.w_k = nn.Linear(config.d_model, config.d_model, bias=False)
        self.w_v = nn.Linear(config.d_model, config.d_model, bias=False)
        self.w_o = nn.Linear(config.d_model, config.d_model, bias=False)

        self.feature_map = self._get_feature_map(config.feature_map)

    def _get_feature_map(self, name: str) -> FeatureMap:
        feature_maps = {
            "elu": ELUFeatureMap(),
        }
        return feature_maps[name]

    def forward(
        self,
        x: Tensor,
        mask: Tensor | None = None
    ) -> Tensor:
        batch_size, seq_len, _ = x.shape

        # Linear projections
        q = self.w_q(x).view(batch_size, seq_len, self.config.n_heads, self.d_k)
        k = self.w_k(x).view(batch_size, seq_len, self.config.n_heads, self.d_k)
        v = self.w_v(x).view(batch_size, seq_len, self.config.n_heads, self.d_k)

        # Apply feature map
        q = self.feature_map(q)
        k = self.feature_map(k)

        # Compute attention in linear time
        # (B, H, L, D) -> (B, H, D, D)
        kv = torch.einsum('bhld,bhle->bhde', k, v)

        # (B, H, L, D) x (B, H, D, D) -> (B, H, L, D)
        out = torch.einsum('bhld,bhde->bhle', q, kv)

        # Normalize
        k_sum = torch.einsum('bhld->bhl', k).unsqueeze(-1)
        out = out / (torch.einsum('bhld,bhl->bhld', q, k_sum) + self.config.eps)

        # Reshape and project
        out = out.transpose(1, 2).contiguous().view(batch_size, seq_len, -1)
        return self.w_o(out)
\end{lstlisting}

\subsubsection{Sparse Attention}

Sparse attention restricts attention to specific patterns, reducing computation.

\begin{lstlisting}[language=Python, caption=Sparse Attention Patterns]
@dataclass(frozen=True, slots=True)
class SparseAttentionPattern:
    """Base class for sparse attention patterns"""
    n_positions: int

class LocalAttentionPattern(SparseAttentionPattern):
    """Attend only to local window"""
    window_size: int

    def get_mask(self) -> Tensor:
        mask = torch.zeros(self.n_positions, self.n_positions)
        for i in range(self.n_positions):
            start = max(0, i - self.window_size // 2)
            end = min(self.n_positions, i + self.window_size // 2 + 1)
            mask[i, start:end] = 1
        return mask.bool()

class StridedAttentionPattern(SparseAttentionPattern):
    """Attend with stride (for hierarchical processing)"""
    stride: int

    def get_mask(self) -> Tensor:
        mask = torch.zeros(self.n_positions, self.n_positions)
        for i in range(self.n_positions):
            # Local attention
            mask[i, max(0, i-1):min(self.n_positions, i+2)] = 1
            # Strided attention
            for j in range(0, self.n_positions, self.stride):
                mask[i, j] = 1
        return mask.bool()

class DilatedAttentionPattern(SparseAttentionPattern):
    """Dilated attention for increasing receptive field"""
    dilation_rates: list[int]

    def get_mask(self) -> Tensor:
        mask = torch.zeros(self.n_positions, self.n_positions)
        for i in range(self.n_positions):
            for rate in self.dilation_rates:
                # Attend to positions at multiples of dilation rate
                j = i - rate
                while j >= 0:
                    mask[i, j] = 1
                    j -= rate
                j = i + rate
                while j < self.n_positions:
                    mask[i, j] = 1
                    j += rate
        return mask.bool()
\end{lstlisting}

\subsubsection{Flash Attention}

Flash Attention uses tiling and recomputation to achieve memory-efficient exact attention:

\begin{lstlisting}[language=Python, caption=Flash Attention Concept]
class FlashAttentionConcept:
    """Conceptual implementation of Flash Attention algorithm"""

    @staticmethod
    def tiled_attention(
        q: Tensor,  # (batch, n_heads, seq_len, d_head)
        k: Tensor,
        v: Tensor,
        block_size: int = 64
    ) -> Tensor:
        """
        Compute attention using tiling to reduce memory usage.
        This is a simplified conceptual version.
        """
        batch, n_heads, seq_len, d_head = q.shape

        # Initialize output
        o = torch.zeros_like(q)
        l = torch.zeros(batch, n_heads, seq_len, 1, device=q.device)

        # Process in blocks
        n_blocks = (seq_len + block_size - 1) // block_size

        for i in range(n_blocks):
            start_i = i * block_size
            end_i = min((i + 1) * block_size, seq_len)

            # Load Q block
            q_block = q[:, :, start_i:end_i]
            o_block = torch.zeros_like(q_block)
            l_block = torch.zeros(
                batch, n_heads, end_i - start_i, 1,
                device=q.device
            )

            # Iterate over K, V blocks
            for j in range(n_blocks):
                start_j = j * block_size
                end_j = min((j + 1) * block_size, seq_len)

                # Load K, V blocks
                k_block = k[:, :, start_j:end_j]
                v_block = v[:, :, start_j:end_j]

                # Compute attention scores for this block
                scores = torch.matmul(
                    q_block, k_block.transpose(-2, -1)
                ) / (d_head ** 0.5)

                # For causal attention, mask future positions
                if i < j:
                    scores.fill_(-float('inf'))
                elif i == j:
                    mask = torch.triu(
                        torch.ones(end_i - start_i, end_j - start_j),
                        diagonal=1
                    ).bool()
                    scores.masked_fill_(mask, -float('inf'))

                # Compute attention weights
                p = F.softmax(scores, dim=-1)

                # Update running statistics
                l_new = torch.exp(scores).sum(dim=-1, keepdim=True)
                l_block = l_block + l_new

                # Accumulate weighted values
                o_block = o_block + torch.matmul(p, v_block)

            # Write back results
            o[:, :, start_i:end_i] = o_block
            l[:, :, start_i:end_i] = l_block

        # Final normalization
        o = o / l
        return o
\end{lstlisting}

\subsection{Attention with External Memory}

\subsubsection{Memory-Augmented Transformers}

Adding external memory banks allows transformers to handle longer contexts:

\begin{lstlisting}[language=Python, caption=Memory-Augmented Transformer]
@dataclass(frozen=True, slots=True)
class MemoryBankConfig:
    memory_size: int
    memory_dim: int
    n_heads: int
    update_method: str = "gru"  # "gru", "linear", "none"

class MemoryBank(nn.Module):
    """External memory bank for transformers"""

    def __init__(self, config: MemoryBankConfig) -> None:
        super().__init__()
        self.config = config

        # Initialize memory
        self.register_buffer(
            "memory",
            torch.randn(config.memory_size, config.memory_dim)
        )

        # Memory update mechanism
        if config.update_method == "gru":
            self.update_gate = nn.GRUCell(
                config.memory_dim,
                config.memory_dim
            )
        elif config.update_method == "linear":
            self.update_linear = nn.Linear(
                config.memory_dim * 2,
                config.memory_dim
            )

    def query(
        self,
        queries: Tensor,  # (batch, seq_len, d_model)
        values: Tensor | None = None
    ) -> tuple[Tensor, Tensor]:
        """Query memory bank and optionally update it"""

        batch_size, seq_len, d_model = queries.shape

        # Expand memory for batch processing
        memory = self.memory.unsqueeze(0).expand(batch_size, -1, -1)

        # Compute attention to memory
        scores = torch.matmul(
            queries, memory.transpose(-2, -1)
        ) / (self.config.memory_dim ** 0.5)

        attn_weights = F.softmax(scores, dim=-1)
        memory_output = torch.matmul(attn_weights, memory)

        # Optionally update memory
        if values is not None and self.config.update_method != "none":
            self._update_memory(attn_weights, values)

        return memory_output, attn_weights

    def _update_memory(
        self,
        attn_weights: Tensor,  # (batch, seq_len, memory_size)
        values: Tensor  # (batch, seq_len, d_model)
    ) -> None:
        """Update memory based on attention patterns"""

        # Aggregate updates across sequence
        write_weights = attn_weights.mean(dim=1)  # (batch, memory_size)
        write_values = torch.matmul(
            write_weights.transpose(-2, -1),
            values.mean(dim=1)
        )  # (memory_size, d_model)

        if self.config.update_method == "gru":
            # GRU-style update
            self.memory = self.update_gate(
                write_values,
                self.memory
            )
        elif self.config.update_method == "linear":
            # Linear interpolation
            combined = torch.cat([self.memory, write_values], dim=-1)
            self.memory = self.update_linear(combined)

class MemoryAugmentedTransformerLayer(nn.Module):
    """Transformer layer with external memory"""

    def __init__(
        self,
        d_model: int,
        n_heads: int,
        d_ff: int,
        memory_config: MemoryBankConfig,
        dropout: float = 0.1
    ) -> None:
        super().__init__()

        # Standard transformer components
        self.self_attn = nn.MultiheadAttention(
            d_model, n_heads, dropout=dropout, batch_first=True
        )
        self.memory_bank = MemoryBank(memory_config)
        self.cross_attn = nn.MultiheadAttention(
            d_model, n_heads, dropout=dropout, batch_first=True
        )

        # Feed-forward network
        self.ff = nn.Sequential(
            nn.Linear(d_model, d_ff),
            nn.ReLU(),
            nn.Dropout(dropout),
            nn.Linear(d_ff, d_model)
        )

        # Layer norms
        self.norm1 = nn.LayerNorm(d_model)
        self.norm2 = nn.LayerNorm(d_model)
        self.norm3 = nn.LayerNorm(d_model)

        self.dropout = nn.Dropout(dropout)

    def forward(
        self,
        x: Tensor,
        update_memory: bool = True
    ) -> Tensor:
        # Self-attention
        attn_out, _ = self.self_attn(x, x, x)
        x = self.norm1(x + self.dropout(attn_out))

        # Memory attention
        memory_out, _ = self.memory_bank.query(
            x,
            x if update_memory else None
        )
        x = self.norm2(x + self.dropout(memory_out))

        # Feed-forward
        ff_out = self.ff(x)
        x = self.norm3(x + self.dropout(ff_out))

        return x
\end{lstlisting}

\subsection{Hierarchical and Recursive Attention}

\subsubsection{Hierarchical Transformers}

Process sequences at multiple scales:

\begin{lstlisting}[language=Python, caption=Hierarchical Transformer Architecture]
@dataclass(frozen=True, slots=True)
class HierarchicalConfig:
    levels: list[int]  # e.g., [4, 16, 64] for 3-level hierarchy
    d_model: int
    n_heads: int
    merge_method: str = "mean"  # "mean", "max", "learned"

class HierarchicalTransformer(nn.Module):
    """Multi-scale hierarchical transformer"""

    def __init__(self, config: HierarchicalConfig) -> None:
        super().__init__()
        self.config = config

        # Transformer for each level
        self.transformers = nn.ModuleList([
            nn.TransformerEncoder(
                nn.TransformerEncoderLayer(
                    d_model=config.d_model,
                    nhead=config.n_heads,
                    batch_first=True
                ),
                num_layers=2
            )
            for _ in config.levels
        ])

        # Merging layers between levels
        if config.merge_method == "learned":
            self.merge_layers = nn.ModuleList([
                nn.Linear(config.d_model * 2, config.d_model)
                for _ in range(len(config.levels) - 1)
            ])

    def forward(self, x: Tensor) -> list[Tensor]:
        """Process input hierarchically"""

        batch_size, seq_len, d_model = x.shape
        outputs = []

        # Process each level
        current_input = x
        for level_idx, level_size in enumerate(self.config.levels):
            # Reshape input for this level
            if seq_len % level_size != 0:
                # Pad if necessary
                pad_len = level_size - (seq_len % level_size)
                current_input = F.pad(
                    current_input,
                    (0, 0, 0, pad_len)
                )

            # Group sequences
            grouped = current_input.view(
                batch_size,
                -1,
                level_size,
                d_model
            )

            # Apply transformer at this level
            level_outputs = []
            for group_idx in range(grouped.shape[1]):
                group = grouped[:, group_idx]
                transformed = self.transformers[level_idx](group)
                level_outputs.append(transformed)

            level_output = torch.stack(level_outputs, dim=1)
            outputs.append(level_output)

            # Prepare input for next level
            if level_idx < len(self.config.levels) - 1:
                current_input = self._merge_groups(
                    level_output,
                    level_idx
                )

        return outputs

    def _merge_groups(
        self,
        grouped_output: Tensor,
        level_idx: int
    ) -> Tensor:
        """Merge groups for next hierarchical level"""

        batch_size, n_groups, group_size, d_model = grouped_output.shape

        if self.config.merge_method == "mean":
            # Average pooling
            merged = grouped_output.mean(dim=2)
        elif self.config.merge_method == "max":
            # Max pooling
            merged = grouped_output.max(dim=2)[0]
        elif self.config.merge_method == "learned":
            # Learned merging with first and last tokens
            first_last = torch.cat([
                grouped_output[:, :, 0],
                grouped_output[:, :, -1]
            ], dim=-1)
            merged = self.merge_layers[level_idx](first_last)

        return merged.view(batch_size, -1, d_model)
\end{lstlisting}

\subsubsection{Tree Transformers}

Process tree-structured data with attention:

\begin{lstlisting}[language=Python, caption=Tree Transformer Implementation]
@dataclass(frozen=True)
class TreeNode:
    """Node in a tree structure"""
    value: Tensor
    children: list[TreeNode]
    parent: TreeNode | None = None

class TreeTransformer(nn.Module):
    """Transformer for tree-structured data"""

    def __init__(
        self,
        d_model: int,
        n_heads: int,
        n_layers: int,
        max_depth: int
    ) -> None:
        super().__init__()

        self.d_model = d_model
        self.n_heads = n_heads
        self.max_depth = max_depth

        # Positional encoding for tree depth
        self.depth_embedding = nn.Embedding(max_depth, d_model)

        # Transformer layers
        self.layers = nn.ModuleList([
            TreeTransformerLayer(d_model, n_heads)
            for _ in range(n_layers)
        ])

    def forward(self, tree: TreeNode) -> TreeNode:
        """Process tree with transformer"""

        # Convert tree to tensor representation
        nodes, edges, depths = self._tree_to_tensors(tree)

        # Add depth embeddings
        depth_emb = self.depth_embedding(depths)
        nodes = nodes + depth_emb

        # Apply transformer layers
        for layer in self.layers:
            nodes = layer(nodes, edges)

        # Convert back to tree
        output_tree = self._tensors_to_tree(nodes, tree)
        return output_tree

    def _tree_to_tensors(
        self,
        tree: TreeNode
    ) -> tuple[Tensor, Tensor, Tensor]:
        """Convert tree to tensor representation"""

        nodes = []
        edges = []
        depths = []

        def traverse(node: TreeNode, depth: int = 0) -> int:
            node_idx = len(nodes)
            nodes.append(node.value)
            depths.append(depth)

            for child in node.children:
                child_idx = traverse(child, depth + 1)
                edges.append((node_idx, child_idx))
                edges.append((child_idx, node_idx))  # Bidirectional

            return node_idx

        traverse(tree)

        # Convert to tensors
        nodes_tensor = torch.stack(nodes)
        depths_tensor = torch.tensor(depths, dtype=torch.long)

        # Create adjacency matrix
        n_nodes = len(nodes)
        adj_matrix = torch.zeros(n_nodes, n_nodes)
        for i, j in edges:
            adj_matrix[i, j] = 1

        return nodes_tensor, adj_matrix, depths_tensor

    def _tensors_to_tree(
        self,
        nodes: Tensor,
        original_tree: TreeNode
    ) -> TreeNode:
        """Convert tensor representation back to tree"""

        idx = 0

        def rebuild(original_node: TreeNode) -> TreeNode:
            nonlocal idx
            new_node = TreeNode(
                value=nodes[idx],
                children=[],
                parent=None
            )
            idx += 1

            for child in original_node.children:
                new_child = rebuild(child)
                new_child.parent = new_node
                new_node.children.append(new_child)

            return new_node

        return rebuild(original_tree)

class TreeTransformerLayer(nn.Module):
    """Single layer of tree transformer"""

    def __init__(self, d_model: int, n_heads: int) -> None:
        super().__init__()

        self.self_attn = nn.MultiheadAttention(
            d_model,
            n_heads,
            batch_first=True
        )
        self.ff = nn.Sequential(
            nn.Linear(d_model, d_model * 4),
            nn.ReLU(),
            nn.Linear(d_model * 4, d_model)
        )
        self.norm1 = nn.LayerNorm(d_model)
        self.norm2 = nn.LayerNorm(d_model)

    def forward(
        self,
        nodes: Tensor,
        adj_matrix: Tensor
    ) -> Tensor:
        """Apply tree-aware attention"""

        # Create attention mask from adjacency matrix
        # Allow attention only between connected nodes
        mask = (adj_matrix == 0).float() * -1e9

        # Self-attention with tree structure
        attn_out, _ = self.self_attn(
            nodes.unsqueeze(0),
            nodes.unsqueeze(0),
            nodes.unsqueeze(0),
            attn_mask=mask.unsqueeze(0)
        )
        nodes = self.norm1(nodes + attn_out.squeeze(0))

        # Feed-forward
        ff_out = self.ff(nodes)
        nodes = self.norm2(nodes + ff_out)

        return nodes
\end{lstlisting}

%==============================================================================
